\chapter{Results}
\label{ch:results}
\noindent
This chapter presents the results from implementing and exercising the proposed
governed ETL architecture. The results are intentionally scoped to the
implemented architecture. They demonstrate whether the architecture can produce
the forms of evidence required by the evaluation framework in Chapter~3. The
formal assessment of those results against the requirements is provided in
Chapter~\ref{ch:evaluation}.

The implemented prototype uses the governed EDGAR pipeline described in
Chapter~5. The pipeline ingests a representative SEC~EDGAR fixture, writes raw
and curated records through the S3 storage zones, validates the curated output
against a data contract, promotes the passing dataset to the governed gold
table layer, and records run evidence in JSON manifests and metadata artifacts.
This setup was sufficient to test the central claim of the architecture: that a
produced dataset version can be reconstructed back to the run, code bundle,
release manifest, infrastructure reference, and change record that produced it.

\section{Implementation Details}
\noindent
The implementation continues the architecture described in Chapter~5 by
materialising the proposed controls across the target orchestration,
infrastructure, metadata, and evidence planes. The pipeline was implemented as a
governed Airflow workflow intended for MWAA, with Terraform defining the AWS
resources, S3 separating the raw, curated, manifest, and gold data zones, and
the metadata layer recording lineage, quality, release, and reconstruction
evidence. The important implementation result is that each produced dataset
version carries enough references to connect the data product back to the
release manifest, change record, code bundle, infrastructure state, quality
result, OpenLineage event, and OpenMetadata records.

The target implementation therefore follows the same split used in the
architecture chapter. The orchestration plane schedules and executes the EDGAR
workflow through Airflow on MWAA. The data plane stores source, intermediate,
manifest, and promoted outputs in the S3-backed zones. The governance plane
uses GitHub Actions, OPA policy checks, release manifests, and artifact
attestation to control change. The evidence plane then persists the run
manifest, audit-chain file, lineage event, catalogue record, and schema
contract so that the produced gold dataset can be reconstructed without relying
on memory of the deployment.

The main implementation hurdle was the cost of repeatedly provisioning the
managed orchestration environment. The MWAA environment took approximately
45~minutes to create from Terraform and approximately 15~minutes to tear down.
That made full rebuilds too slow for normal development, so a Docker-based
Airflow environment was used during local development to allow faster
iteration before deployment to the target architecture.

The second major issue during development was the recent instability of
GitHub. Frequent outages affected the ability to develop and test the system,
with particular frustration caused by the GitHub merge-queue incident. Because
the target architecture depends on GitHub for issues, pull requests, release
controls, ownership, and attestation, these incidents reinforced the need to
retain independent release manifests, hashes, and run manifests rather than
relying only on platform state. These issues create uncertainty in a system
where the foundational principle is trust. Questions about reliance on GitHub
are therefore returned to later in the thesis.

The detailed implementation version and platform reference tables are provided
in Appendix~\ref{chap:implementation-versions} so that the results discussion
can focus on the evidence produced by the architecture rather than the support
stack inventory.

\section{Observed Results}
\noindent
The results in this chapter are presented as a worked example following one
representative governed EDGAR run. The purpose of doing this is to show how the
architecture behaves when an auditor starts from a produced dataset and follows
the evidence chain backward. The example run ingested the representative
SEC~EDGAR fixture, transformed it into the curated contract shape, executed the
quality gate, promoted the passing output to the gold table, and wrote the
surrounding run, lineage, catalogue, and audit-chain evidence.

For this demonstration run, the promoted gold dataset version was
\texttt{9b5fd26a2005c02b}. The value is used throughout the remainder of this
chapter as the example lookup key, not as a claim that the identifier itself is
important outside this run. Repeated execution over the same governed input,
contract, and transformation logic resolved to the same version identifier,
which is the expected behaviour because the version is derived from promoted
content and table identity. The important result was therefore not the volume
of data processed, but that the run produced a stable governed output with
enough surrounding evidence to explain that output without manual log
collation.

Table~\ref{tab:results-artifact-linkage} summarises how the artifacts in this
single worked example point to each other. The table should be read as the
traceback path for the demonstration run rather than as a general inventory of
all possible evidence files.

\begin{table}[H]
  \centering
  \caption{Evidence artifact linkage used for audit reconstruction.}
  \label{tab:results-artifact-linkage}
  \renewcommand{\arraystretch}{1.15}
  \begin{tabularx}{\textwidth}{@{}p{0.25\textwidth}p{0.26\textwidth}X@{}}
    \toprule
    \textbf{Artifact} & \textbf{Pointer field} & \textbf{What it resolves to} \\
    \midrule
    Gold dataset & \texttt{dataset\_version} & Resolves to the audit-chain artifact for \texttt{9b5fd26a2005c02b}. \\
    Audit-chain artifact & \texttt{manifest\_ref} & Resolves to the run manifest that produced the dataset. \\
    Run manifest & \texttt{release\_manifest\_ref} & Resolves to the release manifest that defines the DAG and job files. \\
    Run manifest & \texttt{change\_ref} & Resolves to the governed change record, including the \texttt{EDGAR-1} implementation record. \\
    Run manifest & \texttt{terraform\_state\_ref} & Resolves to the infrastructure state reference for the execution environment. \\
    Run manifest & \texttt{curated\_outputs.contract\_path} & Resolves to the curated contract used by the quality gate. \\
    Run manifest & \texttt{code\_bundle.artifacts} & Resolves to hashed DAG, ingest, transform, quality, and promotion files. \\
    Run manifest & \texttt{metadata\_refs} & Resolves to OpenLineage, OpenMetadata, and audit-chain metadata records. \\
    OpenLineage event & \texttt{inputs} and \texttt{outputs} & Resolves source, raw, curated, manifest, and gold-table lineage endpoints. \\
    OpenMetadata dataset record & \texttt{upstreams} and \texttt{manifest\_ref} & Resolves the dataset back to upstream inputs and the producing run manifest. \\
    \bottomrule
  \end{tabularx}
\end{table}

In this demonstration run, the quality gate showed that the curated dataset
existed, matched the declared contract, contained non-empty records, and had
required identifiers populated. These are deliberately simple controls, but
they demonstrate the runtime enforcement pattern: a dataset is not promoted
because a task finishes; it is promoted because the required checks pass and
their results are recorded.

The generated run manifest is the main evidence object for the example. For the
successful run, it records the run identifier, source reference, raw outputs,
curated output, contract path, row counts, quality result, gold table
coordinates, dataset version, release manifest reference, Terraform state
reference, and change reference. In the complete governed Airflow run, the
manifest also records source-control provenance, code bundle hashes,
attestation fields, reference checks, OpenLineage output,
OpenMetadata-style records, and the audit-chain path. This satisfies the
practical requirement that the output can be explained without manually
reconstructing the run from unrelated logs.

\section{Evidence Schemas}
\noindent
The results are grounded in the structure of the evidence artifacts rather than
only in task success. Each artifact below is materialised as JSON and records a
different part of the governed chain. In the worked example, these schemas
explain what each evidence artifact is expected to carry as the run moves from
ingestion through promotion and then into audit reconstruction.

\begin{table}[H]
  \centering
  \caption{Run manifest schema.}
  \label{tab:schema-run-manifest}
  \renewcommand{\arraystretch}{1.12}
  \begin{tabularx}{\textwidth}{@{}p{0.3\textwidth}X@{}}
    \toprule
    \textbf{Field} & \textbf{Explanation} \\
    \midrule
    \texttt{run\_id} & Identifies the specific governed execution. \\
    \texttt{pipeline\_id} & Identifies the pipeline definition that ran. \\
    \texttt{dataset\_date} & Identifies the dataset partition being produced. \\
    \texttt{source\_ref} & Records the upstream source object or fixture used for ingestion. \\
    \texttt{release\_manifest\_ref} & Links execution to the approved release definition. \\
    \texttt{terraform\_state\_ref} & Links execution to the infrastructure state reference. \\
    \texttt{change\_ref} & Links execution to the governing change record. \\
    \texttt{job\_refs} & Lists the ingest, transform, quality, and promotion jobs used by the run. \\
    \texttt{raw\_outputs} & Records raw payload and normalized raw output locations. \\
    \texttt{curated\_outputs} & Records curated output and contract path. \\
    \texttt{gold\_write\_result} & Records promoted gold table write details. \\
    \texttt{row\_counts} & Stores raw, curated, and promoted count consistency evidence. \\
    \texttt{quality\_result} & Captures the blocking quality-gate result and per-check outcomes. \\
    \texttt{dataset\_version} & Identifies the promoted gold dataset version. \\
    \texttt{reference\_checks} & Records whether release, infrastructure, and change references resolved at runtime. \\
    \texttt{source\_control} & Captures repository, branch, commit, and commit URL evidence. \\
    \texttt{code\_bundle} & Captures release-manifest and job artifact hashes. \\
    \texttt{attestation} & Captures GitHub attestation provider and subject fields. \\
    \texttt{metadata\_refs} & Links to catalogue, lineage, and audit-chain artifacts. \\
    \texttt{status} & Records the final run state. \\
    \texttt{started\_at} & Records when the manifest began tracking the run. \\
    \texttt{completed\_at} & Records when the run evidence was finalised. \\
    \bottomrule
  \end{tabularx}
\end{table}

\begin{table}[H]
  \centering
  \caption{Release manifest schema.}
  \label{tab:schema-release-manifest}
  \renewcommand{\arraystretch}{1.12}
  \begin{tabularx}{\textwidth}{@{}p{0.3\textwidth}X@{}}
    \toprule
    \textbf{Field} & \textbf{Explanation} \\
    \midrule
    \texttt{release\_id} & Names the governed release, such as \texttt{dev-edgar-v1}. \\
    \texttt{pipeline\_id} & Binds the release to \texttt{edgar\_governed\_pipeline}. \\
    \texttt{dag\_bundle\_ref} & Points to the Airflow DAG bundle included in the release. \\
    \texttt{jobs} & Lists the ingest, transform, quality, and promotion job files. \\
    \texttt{source\_control} & Records commit, branch, repository, and commit URL when provided by the release workflow. \\
    \texttt{artifact\_bundle} & Reserves the package path, digest, subject digest, and attestation bundle path. \\
    \texttt{attestation} & Records GitHub attestation provider, run, workflow, URL, and subject digest fields. \\
    \texttt{notes} & Documents the release purpose or environment. \\
    \bottomrule
  \end{tabularx}
\end{table}

\begin{table}[H]
  \centering
  \caption{Change record schema.}
  \label{tab:schema-change-record}
  \renewcommand{\arraystretch}{1.12}
  \begin{tabularx}{\textwidth}{@{}p{0.3\textwidth}X@{}}
    \toprule
    \textbf{Field} & \textbf{Explanation} \\
    \midrule
    \texttt{change\_id} & Names the governing change, for example \texttt{EDGAR-1}. \\
    \texttt{title} & Describes the implementation or rehearsal governed by the change. \\
    \texttt{source} & Identifies the governance source used for the change record. \\
    \texttt{source\_control} & Stores commit, branch, repository, and commit URL fields when available. \\
    \texttt{attestation} & Mirrors the GitHub attestation fields used by release and run evidence. \\
    \bottomrule
  \end{tabularx}
\end{table}

\begin{table}[H]
  \centering
  \caption{Curated contract schema.}
  \label{tab:schema-curated-contract}
  \renewcommand{\arraystretch}{1.12}
  \begin{tabularx}{\textwidth}{@{}p{0.3\textwidth}X@{}}
    \toprule
    \textbf{Field} & \textbf{Explanation} \\
    \midrule
    \texttt{dataset} & Names the contract target, \texttt{edgar\_curated}. \\
    \texttt{description} & Explains the deterministic governed EDGAR filing slice. \\
    \texttt{required\_fields} & Requires \texttt{accession\_number} and \texttt{company\_cik}. \\
    \texttt{columns} & Defines 8 expected curated columns and their string types. \\
    \bottomrule
  \end{tabularx}
\end{table}

\begin{table}[H]
  \centering
  \caption{Audit-chain schema.}
  \label{tab:schema-audit-chain}
  \renewcommand{\arraystretch}{1.12}
  \begin{tabularx}{\textwidth}{@{}p{0.3\textwidth}X@{}}
    \toprule
    \textbf{Field} & \textbf{Explanation} \\
    \midrule
    \texttt{dataset\_version} & Starts reconstruction from the promoted gold dataset version. \\
    \texttt{gold\_table\_uri} & Identifies the promoted gold table location. \\
    \texttt{pipeline\_id} & Identifies the pipeline that produced the dataset. \\
    \texttt{run\_id} & Identifies the producing run. \\
    \texttt{manifest\_ref} & Resolves the dataset to the producing run manifest. \\
    \texttt{source\_ref} & Identifies the source input used by the run. \\
    \texttt{release\_manifest\_ref} & Links the dataset to release evidence. \\
    \texttt{terraform\_state\_ref} & Links the dataset to infrastructure evidence. \\
    \texttt{change\_ref} & Links the dataset to change evidence. \\
    \texttt{source\_control} & Records commit provenance. \\
    \texttt{code\_bundle} & Records release hashes and job hashes. \\
    \texttt{attestation} & Records attestation provider and subject fields. \\
    \texttt{quality\_status} & Records whether the dataset passed the quality gate before promotion. \\
    \bottomrule
  \end{tabularx}
\end{table}

\begin{table}[H]
  \centering
  \caption{OpenLineage event schema.}
  \label{tab:schema-openlineage}
  \renewcommand{\arraystretch}{1.12}
  \begin{tabularx}{\textwidth}{@{}p{0.3\textwidth}X@{}}
    \toprule
    \textbf{Field} & \textbf{Explanation} \\
    \midrule
    \texttt{eventType} & Records the lineage event type. \\
    \texttt{eventTime} & Records the lineage event time. \\
    \texttt{job} & Identifies the pipeline namespace and job name. \\
    \texttt{run} & Identifies the OpenLineage run UUID and governance facets. \\
    \texttt{inputs} & Lists source, raw, and curated upstream datasets. \\
    \texttt{outputs} & Lists the run manifest and promoted S3~Tables dataset. \\
    \texttt{facets.version} & Stores dataset version, content digest, and promoted records path. \\
    \texttt{producer} & Records the producing metadata component. \\
    \texttt{schemaURL} & Records the OpenLineage schema reference. \\
    \bottomrule
  \end{tabularx}
\end{table}

\begin{table}[H]
  \centering
  \caption{OpenMetadata run-record schema.}
  \label{tab:schema-openmetadata-run}
  \renewcommand{\arraystretch}{1.12}
  \begin{tabularx}{\textwidth}{@{}p{0.3\textwidth}X@{}}
    \toprule
    \textbf{Field} & \textbf{Explanation} \\
    \midrule
    \texttt{entity\_type} & Marks the record as a pipeline run. \\
    \texttt{pipeline\_id} & Identifies the producing pipeline. \\
    \texttt{run\_id} & Identifies the producing run. \\
    \texttt{dataset\_date} & Identifies the dataset partition. \\
    \texttt{status} & Records run status. \\
    \texttt{quality\_status} & Records quality-gate status. \\
    \texttt{row\_counts} & Records count consistency. \\
    \texttt{manifest\_ref} & Links the record back to the run manifest. \\
    \texttt{source\_ref} & Links the record to the source input. \\
    \texttt{release\_manifest\_ref} & Preserves the release reference. \\
    \texttt{terraform\_state\_ref} & Preserves the infrastructure reference. \\
    \texttt{change\_ref} & Preserves the change reference. \\
    \texttt{source\_control} & Carries commit provenance into catalogue metadata. \\
    \texttt{code\_bundle} & Carries artifact hashes into catalogue metadata. \\
    \texttt{attestation} & Carries attestation fields into catalogue metadata. \\
    \bottomrule
  \end{tabularx}
\end{table}

\begin{table}[H]
  \centering
  \caption{OpenMetadata dataset-record schema.}
  \label{tab:schema-openmetadata-dataset}
  \renewcommand{\arraystretch}{1.12}
  \begin{tabularx}{\textwidth}{@{}p{0.3\textwidth}X@{}}
    \toprule
    \textbf{Field} & \textbf{Explanation} \\
    \midrule
    \texttt{entity\_type} & Marks the record as a dataset. \\
    \texttt{dataset\_version} & Identifies the promoted dataset version. \\
    \texttt{content\_digest} & Fingerprints the promoted dataset content. \\
    \texttt{gold\_table\_bucket} & Identifies the gold table bucket. \\
    \texttt{gold\_namespace} & Identifies the gold table namespace. \\
    \texttt{gold\_table} & Identifies the gold table name. \\
    \texttt{table\_uri} & Locates the promoted gold table. \\
    \texttt{records\_path} & Locates the materialised records file. \\
    \texttt{write\_result\_path} & Locates the write-result evidence file. \\
    \texttt{row\_count} & Records count evidence. \\
    \texttt{run\_id} & Links the dataset back to the producing run. \\
    \texttt{manifest\_ref} & Links the dataset back to the producing run manifest. \\
    \texttt{source\_curated\_uri} & Records the curated source used for promotion. \\
    \texttt{upstreams} & Preserves upstream lineage inputs. \\
    \texttt{source\_control} & Carries code provenance into the dataset record. \\
    \texttt{attestation} & Carries attestation fields into the dataset record. \\
    \bottomrule
  \end{tabularx}
\end{table}

\section{Audit Reconstruction Result}
\noindent
The strongest result from the worked example is the audit reconstruction path.
Starting only from the gold dataset version
\texttt{9b5fd26a2005c02b}, an auditor can resolve the audit-chain artifact at
\path{metadata/audit-chain/9b5fd26a2005c02b.json}. Within this example, that
artifact identifies the pipeline \texttt{edgar\_governed\_pipeline}, the
Airflow run that produced the output, the source fixture, the gold table URI
\path{s3tables://gold/edgar/filings}, the release manifest reference, the
Terraform state reference, the source-control commit, and the code bundle
hashes for the DAG and job scripts.

The reconstruction then moves from the audit-chain artifact to the run
manifest. The manifest records the raw payload and normalized raw record
locations, the curated output location, the contract path, the quality result,
and the promoted gold write result. The promoted write result records the
table URI, dataset date, content digest, source curated URI, write-result path,
and dataset version. This means the output can be traced backward from the
gold table to the curated file, and then backward again to the raw source
payload.

The same manifest also reconstructs the control path for the example run. The
\texttt{reference\_checks} block records that
\path{release-manifests/dev-edgar-v1.json},
\path{terraform-state/aws-target.json}, and
\path{changes/EDGAR-1.json} resolved successfully. The
\texttt{code\_bundle} block records the release-manifest SHA-256 hash and
individual hashes for \path{pipelines/edgar_governed_pipeline.py},
\path{jobs/ingest/edgar_ingest.py},
\path{jobs/transform/edgar_transform.py},
\path{jobs/quality/curated_quality_gate.py}, and
\path{jobs/promote/promote_to_s3_table.py}. These hashes are important
because they let the reconstruction identify the code that was intended to run,
not only the path names of the files.

The schema evidence is also part of the reconstruction. The curated output
points to \path{contracts/edgar_curated_contract.json}, which declares the
expected columns and requires \texttt{accession\_number} and
\texttt{company\_cik}. The quality result then records that the curated file
matched this contract and that required identifiers were present. In this
sense, the schema is not separate documentation. It is executable evidence
used by the quality gate and retained in the manifest.

Finally, metadata publication connects the example run evidence file to the
broader catalogue and lineage model. The \texttt{metadata\_refs} block points to the
OpenLineage completion event, the OpenMetadata-compatible run record, the
OpenMetadata-compatible dataset record, and the audit-chain path. It also
records the pipeline FQN, storage service name, raw, curated, and gold
container FQNs, and lineage edge count. The worked example therefore has two
directions of proof: the manifest can reconstruct the run from the dataset,
and the metadata records can expose the same run through lineage and catalogue
surfaces.

This result directly addresses the weakness observed in the baseline pipeline.
In the baseline, a completed dataset could be associated with logs and output
paths, but there was no single evidence object that bound the output to its
change intent, runtime code, infrastructure reference, and quality decision. In
the proposed implementation, those references are persisted together. The
result is not merely more logging; it is a structured reconstruction path from
dataset version back to operational authority.

\section{Summary}
\noindent
The results show that the proposed architecture can produce a working evidence
chain for a governed ETL run. In the representative worked example, the
prototype generated a structured run manifest, 4 quality checks, 5 hashed code
artifacts, 3 classes of metadata evidence, 6 lineage edges, and a stable
dataset version that links the gold dataset back to code, infrastructure, and
change references.

These results are reported from the representative EDGAR evaluation run and the
development harness used to exercise the target design. That scope is
intentional because the purpose of the results is to test the evidence chain
from release to runtime output. The prototype demonstrates the central
architectural claim: auditability improves when governance evidence is captured
as part of the delivery and runtime workflow, rather than reconstructed manually
after a dataset has already been produced.
